problem:  
Let \(M\) be a smooth, compact manifold carrying an equiregular, bracket-generating distribution \(\mathcal{D}\subset TM\) of step \(2\) endowed with a sub-Riemannian metric \(g_{\mathcal{D}}\) and corresponding Popp measure \(\mu_{\mathrm{Popp}}\).  
For each \(\lambda>0\), consider the Riemannian metric  
\[
g_\lambda = g_{\mathcal{D}} \oplus \lambda^2 g_{\mathcal{D}^\perp}},
\]
where \(\mathcal{D}^\perp\) is a smooth complement chosen orthonormally with respect to a fixed background Riemannian metric.  
Assume that all \((M,g_\lambda)\) satisfy uniform two-sided sectional curvature bounds \(-K\leq \mathrm{Sec}_{g_\lambda}\leq K\) independent of \(\lambda\).  

Under these hypotheses, the sequence \((M,g_\lambda,\mathrm{vol}_{g_\lambda})\) converges (after normalization) in the measured Gromov–Hausdorff sense to a metric-measure space \((M,d_{\mathrm{CC}},\mu_{\mathrm{Popp}})\), where \(d_{\mathrm{CC}}\) is the Carnot–Carathéodory distance.  

**Question:** Determine whether there exists a universal constant \(c>0\) such that for every such family,
\[
(M,d_{\mathrm{CC}},\mu_{\mathrm{Popp}}) \text{ satisfies a generalized curvature-dimension inequality }
CD_{\text{BG}}(cK,N)
\]
in the sense of Baudoin–Garofalo, for some finite \(N\) depending only on \(\dim M\) and \(\mathrm{rank}(\mathcal{D})\).  
Equivalently, can uniform Riemannian sectional curvature control of the approximating sequence guarantee a *nontrivial synthetic curvature lower bound* in the sub-Riemannian limit, measured by the horizontal Laplacian and Popp measure?

If this fails in general, characterize the class of distributions (e.g. contact, quaternionic contact, or fat step-2) for which such a limiting \(CD_{\text{BG}}\)-type curvature bound necessarily exists.  

Why is it a "good" problem:  
This problem formulates a precise bridge between smooth Riemannian curvature bounds and synthetic curvature-dimension conditions in sub-Riemannian geometry within a rigorously defined limit framework. It refines the vague implication of uniform curvature control into a concrete test of analytic stability for curvature-dimension inequalities under singular metric collapse.  
The problem involves tools from several active research domains—collapsing Riemannian geometry, sub-Riemannian analysis, and metric-measure geometry—and confronts the fundamental question of *how classical curvature control degenerates to the synthetic curvature of horizontal Laplacians*. It is unknown whether such stability of curvature-dimension inequalities survives the sub-Riemannian limit even for contact manifolds, and a positive (or even partial) resolution would constitute a major step toward a unified curvature theory extending beyond the Riemannian world.  
Its statement is conceptually simple—compare curvature bounds before and after the sub-Riemannian limit—but the required insight spans geometry, analysis, and metric convergence theory, exemplifying the depth and interdisciplinary nature of a truly good problem.